% Motivation Letter Template for Arina
% Claude fills this template using (in priority order):
%   1. pi_data/<slug>/summary.md   ← READ THIS FIRST — exact quotes, findings, overlaps
%   2. arina_profile.md            ← Arina's research arc, results, skills
%   3. pi_database.md              ← PI contact details, methods vocabulary
%   4. motivation_tips.md          ← writing principles
% Save output as: outputs/motivation_letters/motivation_<slug>.tex
% Compile: pdflatex motivation_<slug>.tex
%
% RULES FOR CLAUDE WHEN FILLING THIS TEMPLATE:
% ──────────────────────────────────────────────────────────────────────────────
% BEFORE WRITING:
%   Read pi_data/<slug>/summary.md.
%   Copy the KEY FINDINGS section into your working memory for P1.
%   Copy the METHODS QUOTES section into your working memory for P3.
%   Copy the OPEN QUESTIONS section into your working memory for P4.
%   Copy the EXACT ANTIBODIES / INSTRUMENTS table into your working memory for P3.
%
% P1  (3–4 sentences) — Hook: cite a SPECIFIC result from KEY FINDINGS.
%   NOT "I am writing to express my interest."
%   NOT "Your lab's work on microglia is fascinating."
%   YES: "Your 2024 Nature Neuroscience paper showed that xenografted human
%         microglia adopt a distinct HLA-driven antigen-presentation state in
%         response to Aβ oligomers (a response absent in mouse microglia), which
%         suggests that therapeutic targeting will require human-specific models."
%
% P2  (4–5 sentences) — Research arc.
%   IGIB wet-lab → Chalmers GRN → this PI as the natural next step.
%   Name one concrete result from arina_profile.md Research Data Notes.
%
% P3  (4–5 sentences) — Technique fit.
%   Use METHODS QUOTES from summary.md. Match exact antibody names, instrument
%   model numbers, and software from the PI's own paper.
%   YES: "The same antibody panel your lab uses on Aβ-seeded sections (IBA1,
%         Wako 019-19741, 1:500; P2RY12, Sigma HPA014518, 1:200) is the approach
%         I optimised at IGIB for Sox2/Tbr2/NeuN on E14 mouse brain."
%
% P4  (3–4 sentences) — Forward-looking.
%   Seed from OPEN QUESTIONS in summary.md.
%   Connect to Arina's Chalmers finding (SPP1+/TREM2+ hubs at Braak 4–5).
%
% P5  (2 sentences) — Close with specific ask.
%   State: PhD position / 20-minute call / rotation. Attach CV and slides.
%   NOT "at your earliest convenience."
%
% FORMATTING RULES (absolute, no exceptions):
%   NO EM DASHES ( — ) anywhere in the output. Replace with comma, colon,
%   semicolon, or parentheses. Search your output for " — " before saving.
%
% TONE RULES (European academic):
%   - Direct and collegial, not formal or stiff
%   - No flattery ("groundbreaking", "prestigious", "world-class", "thrilled")
%   - No hedging ("I believe I might", "I hope to be able to")
%   - Closing line: "Best regards," — not "Yours sincerely," or "Respectfully,"
%   - For kim (DANDRITE Denmark): salutation is "Dear Thomas,"
%   - For all others first contact: "Dear Prof. [Last Name],"
%     except Netherlands/Belgium where "Dear [First Name]," is also fine
%
% WORD COUNT: 400–500 words. Never more than 1.5 pages.
% Apply ALL rules from motivation_tips.md.
% ─────────────────────────────────────────────────────────────────────────────

\documentclass[11pt,a4paper]{letter}

\usepackage[T1]{fontenc}
\usepackage[utf8]{inputenc}
\usepackage{tgheros}
\renewcommand*\familydefault{\sfdefault}
\usepackage[top=2.5cm, bottom=2.5cm, left=2.8cm, right=2.8cm]{geometry}
\usepackage[hidelinks]{hyperref}
\usepackage{microtype}
\usepackage{setspace}

% Colors
\usepackage{xcolor}
\definecolor{accent}{rgb}{0.12,0.31,0.47}
\definecolor{textgrey}{gray}{0.15}

\color{textgrey}
\setstretch{1.15}

% Remove default letter formatting clutter
\makeatletter
\let\@texttop\relax
\makeatother

\address{
  \textbf{\color{accent}Arina}\\
  arinazubair01@gmail.com\\
  +91 7906615993\\
  \href{https://www.linkedin.com/in/arina-430a2124b/}{linkedin.com/in/arina-430a2124b}
}

\date{\today}

\begin{document}

\begin{letter}{
  \textbf{[PI FULL NAME AND TITLE]}\\
  [Role]\\
  [Institution]\\
  [City, Country]
}

\opening{Dear [Prof./Dr.] [Last Name],}

% ── PARAGRAPH 1: Hook + specific paper finding (3–4 sentences) ──────────────
% Open with the scientific problem or a finding from the PI's key paper.
% Do NOT start with "I am writing to express my interest."
% Example opener: "Your 2019 Nature Neuroscience study demonstrated that human
% microglia adopt a complexity of transcriptional states (including a CRM
% signature absent in mouse models) that may fundamentally determine
% vulnerability to Alzheimer's pathology."
[PARAGRAPH 1]

% ── PARAGRAPH 2: Research arc (4–5 sentences) ────────────────────────────────
% Tell the story: IGIB wet-lab → Chalmers GRN → why this PI is the next step.
% Name specific techniques and one concrete observation or result.
% Example: "During my time as a Project Assistant at CSIR-IGIB under Dr. Arpan
% Parichha, I established a Sox2/Tbr2/NeuN three-marker immunofluorescence
% panel on embryonic mouse cryosections..."
[PARAGRAPH 2]

% ── PARAGRAPH 3: Specific fit with this lab (4–5 sentences) ──────────────────
% Name 2–3 techniques this lab uses that Arina has performed.
% Use the PI's vocabulary from pi_database.md.
% Example: "Your group's multi-channel IF panels (IBA1, P2RY12, and TREM2 with
% quantitative Fiji-based morphometry) mirror the approach I developed at
% CSIR-IGIB..."
[PARAGRAPH 3]

% ── PARAGRAPH 4: Forward-looking (3–4 sentences) ─────────────────────────────
% Propose a rough research direction or question that fits the lab's agenda.
% Connect Arina's GRN results to the PI's methods where possible.
% This is a conversation starter, not a contract.
[PARAGRAPH 4]

% ── PARAGRAPH 5: Close (2 sentences) ─────────────────────────────────────────
% State what you are asking for. Do not use "at your earliest convenience."
[PARAGRAPH 5]

\closing{Sincerely,}

\vspace{-6pt}
\textbf{\color{accent}Arina}\\[4pt]
\small
Attached: CV \quad|\quad
\href{[IGIB slides link]}{IGIB Neurodevelopment Slides} \quad|\quad
\href{[Chalmers slides link]}{AD GRN Slides}

\end{letter}
\end{document}
